\documentclass[10pt,a4paper]{amsart}
\usepackage{amsmath,amssymb,amsthm,graphicx}
\usepackage{verbatim}
\newenvironment{reference}{\comment}{\endcomment}
%\newenvironment{reference}{}{}
\newenvironment{slogan}{\comment}{\endcomment}
\newenvironment{history}{\comment}{\endcomment}

% For commutative diagrams we use Xy-pic
\usepackage[all]{xy}

% We use 2cell for 2-commutative diagrams.
\xyoption{2cell}
% Theorem environments.
%
\theoremstyle{plain}
\newtheorem{theorem}[subsection]{Theorem}
\newtheorem{proposition}[subsection]{Proposition}
\newtheorem{lemma}[subsection]{Lemma}

\theoremstyle{definition}
\newtheorem{definition}[subsection]{Definition}
\newtheorem{example}[subsection]{Example}
\newtheorem{exercise}[subsection]{Exercise}
\newtheorem{situation}[subsection]{Situation}

\theoremstyle{remark}
\newtheorem{remark}[subsection]{Remark}
\newtheorem{remarks}[subsection]{Remarks}

\numberwithin{equation}{subsection}

% Macros
%
\def\lim{\mathop{\mathrm{lim}}\nolimits}
\def\colim{\mathop{\mathrm{colim}}\nolimits}
\def\Spec{\mathop{\mathrm{Spec}}}
\def\Hom{\mathop{\mathrm{Hom}}\nolimits}
\def\Ext{\mathop{\mathrm{Ext}}\nolimits}
\def\SheafHom{\mathop{\mathcal{H}\!\mathit{om}}\nolimits}
\def\SheafExt{\mathop{\mathcal{E}\!\mathit{xt}}\nolimits}
\def\Sch{\mathit{Sch}}
\def\Mor{\mathop{\mathrm{Mor}}\nolimits}
\def\Ob{\mathop{\mathrm{Ob}}\nolimits}
\def\Sh{\mathop{\mathit{Sh}}\nolimits}
\def\NL{\mathop{N\!L}\nolimits}
\def\CH{\mathop{\mathrm{CH}}\nolimits}
\def\proetale{{pro\text{-}\acute{e}tale}}
\def\etale{{\acute{e}tale}}
\def\QCoh{\mathit{QCoh}}
\def\Ker{\mathop{\mathrm{Ker}}}
\def\Im{\mathop{\mathrm{Im}}}
\def\Coker{\mathop{\mathrm{Coker}}}
\def\Coim{\mathop{\mathrm{Coim}}}

% Boxtimes
%
\DeclareMathSymbol{\boxtimes}{\mathbin}{AMSa}{"02}

%
% Macros for moduli stacks/spaces
%
\def\QCohstack{\mathcal{QC}\!\mathit{oh}}
\def\Cohstack{\mathcal{C}\!\mathit{oh}}
\def\Spacesstack{\mathcal{S}\!\mathit{paces}}
\def\Quotfunctor{\mathrm{Quot}}
\def\Hilbfunctor{\mathrm{Hilb}}
\def\Curvesstack{\mathcal{C}\!\mathit{urves}}
\def\Polarizedstack{\mathcal{P}\!\mathit{olarized}}
\def\Complexesstack{\mathcal{C}\!\mathit{omplexes}}
% \Pic is the operator that assigns to X its picard group, usage \Pic(X)
% \Picardstack_{X/B} denotes the Picard stack of X over B
% \Picardfunctor_{X/B} denotes the Picard functor of X over B
\def\Pic{\mathop{\mathrm{Pic}}\nolimits}
\def\Picardstack{\mathcal{P}\!\mathit{ic}}
\def\Picardfunctor{\mathrm{Pic}}
\def\Deformationcategory{\mathcal{D}\!\mathit{ef}}
\usepackage{cleveref}
\crefname{originallemma}{original lemma}{original lemmas}
\Crefname{originallemma}{Original lemma}{Original lemmas}
\makeatletter
\newlabel{formal-spaces-lemma-reduction-completion}{{0GB9}{source}{}{sourcepointer-1}{}}
\newlabel{formal-spaces-lemma-reduction-completion@cref}{{[originallemma][0][]0GB9}{[1][0][]source}{}{}{}}
\makeatother
\crefname{originalsection}{original section}{original sections}
\Crefname{originalsection}{Original section}{Original sections}
\makeatletter
\newlabel{formal-spaces-section-quasi-compact}{{0AJ8}{source}{}{sourcepointer-2}{}}
\newlabel{formal-spaces-section-quasi-compact@cref}{{[originalsection][0][]0AJ8}{[1][0][]source}{}{}{}}
\makeatother
\crefname{originallemma}{original lemma}{original lemmas}
\Crefname{originallemma}{Original lemma}{Original lemmas}
\makeatletter
\newlabel{spaces-more-morphisms-lemma-thicken-property-morphisms}{{09ZY}{source}{}{sourcepointer-3}{}}
\newlabel{spaces-more-morphisms-lemma-thicken-property-morphisms@cref}{{[originallemma][0][]09ZY}{[1][0][]source}{}{}{}}
\makeatother
\crefname{originallemma}{original lemma}{original lemmas}
\Crefname{originallemma}{Original lemma}{Original lemmas}
\makeatletter
\newlabel{spaces-morphisms-lemma-base-change-quasi-compact}{{03HF}{source}{}{sourcepointer-4}{}}
\newlabel{spaces-morphisms-lemma-base-change-quasi-compact@cref}{{[originallemma][0][]03HF}{[1][0][]source}{}{}{}}
\makeatother
\title{Quasi-compact f gives quasi-compact formal completion}
\author{The Stacks Project Authors}
\date{stacks-project:restricted.tex:local-snapshot-20261012; GNU FDL 1.2 or later}
\begin{document}
\maketitle
\noindent\textit{Editorial note.} All delivered mathematical text is the author's own, copied from the pinned Stacks Project chapter file restricted.tex; no statement or proof step was rewritten, reordered or removed.\par\medskip
\noindent\textit{Editorial note.} The counted proof is the author's complete argument: the proof identifies the reduction of the formal completion morphism with the map between reduced induced subspaces, observes it factors as a thickening after a quasi-compact base change, and concludes with the thickenings lemma\par\medskip
\noindent\textit{Editorial note.} The label lemma-quasi-compact-gives-quasi-compact names the counted statement, and the proof environment immediately following it in the source is delivered verbatim as its proof.\par\medskip
\noindent\textit{Editorial note.} The counted statement is carried by the author's own lemma environment in the chapter source.\par\medskip
\noindent\textit{Allowed standard prerequisites (editorial).}\begin{itemize}
\item quasi-compact morphisms of formal algebraic spaces and their reduced models (Formal Spaces, Section formal-spaces-section-quasi-compact) with the identification of the reduction of a completion with the reduced induced subspace on T (Formal Spaces, Lemma formal-spaces-lemma-reduction-completion)
\item quasi-compactness stable under base change (Morphisms of Spaces, Lemma spaces-morphisms-lemma-base-change-quasi-compact)
\item thickening property morphisms: a thickening composed with a quasi-compact morphism stays quasi-compact (More on Morphisms of Spaces, Lemma spaces-more-morphisms-lemma-thicken-property-morphisms)
\end{itemize}\medskip
\ifdefined\hypertarget\hypertarget{sourcepointer-1}{}\fi
\noindent\textit{Source reference.} Original lemma 0GB9: sources/stacks-chow/tags.txt, source lines 20364--20364. Prerequisite pointer invoked at proof line(s) 6077; it names a standard fact of a sibling chapter that the counted argument uses, and every step of the counted proof built on it is retained verbatim in the delivered body.\par\medskip
\ifdefined\hypertarget\hypertarget{sourcepointer-2}{}\fi
\noindent\textit{Source reference.} Original section 0AJ8: sources/stacks-chow/tags.txt, source lines 13291--13291. Prerequisite pointer invoked at proof line(s) 6074; it names a standard fact of a sibling chapter that the counted argument uses, and every step of the counted proof built on it is retained verbatim in the delivered body.\par\medskip
\ifdefined\hypertarget\hypertarget{sourcepointer-3}{}\fi
\noindent\textit{Source reference.} Original lemma 09ZY: sources/stacks-chow/tags.txt, source lines 12616--12616. Prerequisite pointer invoked at proof line(s) 6087; it names a standard fact of a sibling chapter that the counted argument uses, and every step of the counted proof built on it is retained verbatim in the delivered body.\par\medskip
\ifdefined\hypertarget\hypertarget{sourcepointer-4}{}\fi
\noindent\textit{Source reference.} Original lemma 03HF: sources/stacks-chow/tags.txt, source lines 4565--4565. Prerequisite pointer invoked at proof line(s) 6084; it names a standard fact of a sibling chapter that the counted argument uses, and every step of the counted proof built on it is retained verbatim in the delivered body.\par\medskip
\begin{lemma}
\label{lemma-quasi-compact-gives-quasi-compact}
In the situation above. If $f$ is quasi-compact, then
$f_{/T}$ is quasi-compact.
\end{lemma}

\begin{proof}
(Quasi-compact morphisms of formal algebraic spaces are discussed in
Formal Spaces, Section \ref{formal-spaces-section-quasi-compact}.)
We have to show that $(X'_{/T'})_{red} \to (X_{/T})_{red}$ is a quasi-compact
morphism of algebraic spaces. By
Formal Spaces, Lemma \ref{formal-spaces-lemma-reduction-completion}
this is the morphism $Z' \to Z$ where $Z' \subset X'$, resp.\ $Z \subset X$
is the reduced induced algebraic space structure on $T'$, resp.\ $T$.
It follows that $Z' \to f^{-1}Z = Z \times_X X'$ is a thickening (a closed
immersion defining an isomorphism on underlying topological spaces).
Since $Z \times_X X' \to Z$ is quasi-compact as a base change of $f$
(Morphisms of Spaces, Lemma
\ref{spaces-morphisms-lemma-base-change-quasi-compact})
we conclude that $Z' \to Z$ is too by
More on Morphisms of Spaces, Lemma
\ref{spaces-more-morphisms-lemma-thicken-property-morphisms}.
\end{proof}

\end{document}
